%% Section 04 — Program Design / Research Methodology
\chapter{Program Design}
\label{ch:program}

\section{Research Questions}
\label{sec:rqs}

\begin{rqbox}
  \begin{enumerate}[label=\textbf{RQ\arabic*.}, leftmargin=*, itemsep=6pt]
    \item How do \textbf{affinity dynamics} — patterns of positive alignment and
          coalition-building — shape the emergence and stability of consensus
          within the Antarctic Treaty System?
    \item How do \textbf{aversion dynamics} — veto behaviour, strategic defection,
          and antagonistic positioning — threaten or reconfigure treaty equilibria?
    \item What computational simulation framework can reproduce empirically observed
          ATS decision outcomes, and what generalisable principles does it reveal
          for multilateral treaty design?
  \end{enumerate}
\end{rqbox}

\section{Theoretical Framework}
\label{sec:theory}

The project is grounded in three complementary theoretical traditions:

\begin{enumerate}
  \item \textbf{Non-cooperative game theory} — models individual state incentives and
        deviation conditions \parencite{Nash1951, Osborne1994}.
  \item \textbf{Opinion dynamics / consensus models} — captures iterative belief
        convergence among parties \parencite{DeGroot1974}.
  \item \textbf{Agent-based simulation} — provides a computational substrate for
        running counterfactual treaty scenarios \parencite{Epstein1996}.
\end{enumerate}

\section{Research Design Overview}
\label{sec:design}

The project is structured into three interconnected papers, each targeting one research
question and one modelling contribution.

\begin{table}[htbp]
  \centering
  \caption{Paper structure and milestones}
  \label{tab:papers}
  \begin{tabular}{clll}
    \toprule
    \textbf{Paper} & \textbf{Focus} & \textbf{Method} & \textbf{Target venue} \\
    \midrule
    1 & Affinity dynamics     & Network analysis, coalitional GT    & TBD \\
    2 & Aversion dynamics     & Veto-player model, ABM              & TBD \\
    3 & Simulation framework  & Agent-based model, RL               & TBD \\
    \bottomrule
  \end{tabular}
\end{table}

\section{Paper 1 — Affinity Dynamics in Multilateral Treaties}
\label{sec:paper1}

% TODO: ~500 words.
%       - Formal definition of affinity between treaty parties
%       - Data sources (ATS meeting records, voting data)
%       - Network construction and community detection
%       - Expected contribution: first systematic map of affinity coalitions in ATS

\section{Paper 2 — Aversion Dynamics and Consensus Breakdown}
\label{sec:paper2}

% TODO: ~500 words.
%       - Formal definition of aversion / veto propensity
%       - Game-theoretic model of strategic blocking
%       - Calibration to historical ATS consensus failures
%       - Expected contribution: predictive model of consensus breakdown conditions

\section{Paper 3 — Integrated Simulation Framework}
\label{sec:paper3}

% TODO: ~500 words.
%       - Architecture of the agent-based model (state agents, treaty rules, environment)
%       - RL-trained agents for adaptive strategy
%       - Validation against historical ATS outcomes
%       - Generalisability tests on other treaty systems (climate, fisheries)

\section{Data Sources and Ethics}
\label{sec:data}

The primary data sources for this project are:

\begin{itemize}
  \item Antarctic Treaty Consultative Meeting (ATCM) final reports and working papers
        (publicly available at \url{https://www.ats.aq/}).
  \item State-level foreign policy position datasets (e.g., UN voting records).
  \item Published network and voting data from the international relations literature.
\end{itemize}

All data are publicly available institutional records; no human-subjects research is
involved and no ethics clearance is required beyond standard QUT research integrity
protocols.

\section{Timeline}
\label{sec:timeline}

\begin{longtable}{lll}
  \caption{Indicative project timeline} \label{tab:timeline} \\
  \toprule
  \textbf{Phase} & \textbf{Activity} & \textbf{Target Date} \\
  \midrule
  \endfirsthead
  \toprule
  \textbf{Phase} & \textbf{Activity} & \textbf{Target Date} \\
  \midrule
  \endhead
  \midrule
  \multicolumn{3}{r}{\small Continued on next page} \\
  \endfoot
  \bottomrule
  \endlastfoot
  1  & Literature review \& theoretical framework      & Month 6  \\
  2  & Paper 1: data collection \& analysis            & Month 12 \\
  3  & Paper 1: writing \& submission                  & Month 15 \\
  4  & Paper 2: model development \& calibration       & Month 18 \\
  5  & Paper 2: writing \& submission                  & Month 21 \\
  6  & Paper 3: simulation framework development       & Month 27 \\
  7  & Paper 3: writing \& submission                  & Month 30 \\
  8  & Thesis integration \& write-up                  & Month 36 \\
  9  & Thesis submission                               & Month 38 \\
\end{longtable}

\section{Expected Contributions}
\label{sec:contributions}

\begin{enumerate}
  \item A formal computational model of consensus dynamics in the ATS, the first of
        its kind at this level of institutional specificity.
  \item A transferable agent-based simulation framework applicable to any
        consensus-based multilateral treaty.
  \item New theoretical constructs (affinity/aversion indices) operationalising
        previously qualitative political science concepts.
  \item Policy-relevant insights into treaty design features that improve stability
        under strategic pressure.
\end{enumerate}
